%%=============================================================================
%% Inleiding
%%=============================================================================

\chapter{\IfLanguageName{dutch}{Inleiding}{Introduction}}%
\label{ch:inleiding}

Softwareprojecten binnen bedrijven en organisaties maken gebruik van talrijke externe bibliotheken, frameworks en tools, de zogenaamde \emph{dependencies}. Deze dependencies worden voortdurend bijgewerkt om nieuwe functionaliteiten, verbeterde prestaties en vooral beveiligingspatches aan te bieden \textcite {Pashchenko2018}. In de praktijk blijkt het voor ontwikkelteams echter moeilijk om alle updates systematisch op te volgen en grondig te evalueren, zeker wanneer een bedrijf meerdere projecten en repositories beheert.

\section{\IfLanguageName{dutch}{Probleemstelling}{Problem Statement}}%
\label{sec:probleemstelling}

Het manueel opvolgen van dependency-updates vormt een groot probleem in softwareontwikkeling. Ontwikkelaars moeten regelmatig changelogs doorzoeken, compatibiliteit controleren en inschatten welke updates veilig kunnen worden doorgevoerd, wat in de praktijk veel tijd vraagt en vaak wordt uitgesteld \textcite {Pashchenko2018,Behutiye2024}. Onderzoek toont aan dat veel softwareprojecten hun dependencies niet systematisch bijhouden, waardoor een aanzienlijk deel verouderd raakt \textcite {Pashchenko2018}. stellen bijvoorbeeld dat een belangrijk percentage kwetsbare dependencies relatief eenvoudig opgelost kan worden door een update, maar dat dit in de praktijk niet gebeurt wegens tijdsdruk en gebrek aan overzicht.
Wanneer het updateproces niet systematisch gebeurt, stapelen onderhoudstaken zich op, waardoor het steeds moeilijker en tijdrovender wordt om de software up-to-date te houden. Dit fenomeen staat bekend als \emph{technische schuld} (technical debt): de achterstand die ontstaat doordat noodzakelijke verbeteringen of updates te lang worden uitgesteld, wat leidt tot hogere onderhoudskosten en complexiteit op lange termijn \textcite {Behutiye2024,Ruiz2024}. Empirisch onderzoek bevestigt dat technische schuld binnen bedrijven reële negatieve gevolgen heeft voor productiviteit en softwarekwaliteit \textcite {Chalmers2020}.
Om deze last te verlichten, zetten veel organisaties al geautomatiseerde dependencybots in, zoals Dependabot\footnote{Dependabot is de ingebouwde GitHub-tool die kwetsbare of verouderde dependencies detecteert en automatisch update-pull-requests aanmaakt.\url{https://github.com/dependabot}} of Renovate.\footnote{Renovate is een open-source tool die automatisch update-pull-requests voor dependencies aanmaakt.\url{https://docs.renovatebot.com} } Deze tools maken wel automatisch pull requests aan zodra er nieuwe versies beschikbaar zijn,\textcite {Erlenhov2022} maar geven doorgaans weinig context over de inhoud en impact van de wijziging. Ontwikkelaars verliezen daardoor nog steeds tijd aan het interpreteren van changelogs, het inschatten van risico's en het beslissen of een update al dan niet moet worden doorgevoerd \textcite {Erlenhov2022}.
Deze bachelorproef richt zich daarom op het ontwikkelen van een \emph{aparte webapplicatie} (een dashboard) die integreert met GitHub-repositories van Springbok. Per project houdt deze applicatie de gebruikte dependencies bij, controleert via publieke registries (zoals npm\footnote{npm is het standaard pakketbeheerplatform voor het JavaScript-ecosysteem.\url{https://www.npmjs.com}}) of er nieuwe versies beschikbaar zijn en verzamelt de relevante changelogs en release notes. Een AI-agent, die in dit onderzoek wordt ontworpen op basis van bestaande AI-modellen maar zonder voorafgebouwd agentframework, analyseert deze informatie, genereert begrijpelijke samenvattingen en beoordeelt in welke mate een update cruciaal is (bijvoorbeeld om veiligheidsredenen) of eerder optioneel.
De AI-agent geeft het kernteam via het dashboard inzicht in vragen zoals: welke dependencies zijn verouderd, wat is er precies veranderd in de nieuwe versie, zijn er bekende beveiligingsrisico's als niet wordt geüpdatet, en lijkt de update technisch laag- of hoogrisicovol.

\section{\IfLanguageName{dutch}{Onderzoeksvraag}{Research question}}%
\label{sec:onderzoeksvraag}

Om het probleem van handmatig dependencybeheer en beperkte context bij updates op te lossen, wordt de volgende hoofdonderzoeksvraag geformuleerd:

\begin{quote}
  Hoe kan een AI-gedreven agent worden ingezet om dependency-updates uit bestaande tools voor automatisch dependencybeheer te analyseren en te beheren, zodat het kernteam dat verantwoordelijk is voor platform- en dependencybeheer efficiënter en met minder handmatig werk dependency-updates kan verwerken?
\end{quote}

Hieruit vloeien de volgende deelvragen voort.

\subsubsection{Deelvragen voor het beter begrijpen van het probleem}
\label{sec:Deelvragen_voor_het_beter_begrijpen_van_het_probleem}

\begin{enumerate}
  \item Wat zijn de belangrijkste uitdagingen en risico's bij het huidige, overwegend manuele dependency management voor het kernteam dat verantwoordelijk is voor platform- en dependencybeheer?
  \item Welke bestaande tools en oplossingen voor automatisch dependencybeheer worden vandaag gebruikt, en welke sterke punten en beperkingen hebben deze oplossingen in de context van Springbok?
  \item Hoe leidt ophopende technische schuld op het vlak van dependency-updates ertoe dat er steeds meer tijd en inspanning nodig zijn voor het controleren en bijwerken van dependencies binnen Springbok?
\end{enumerate}

\subsubsection{Deelvragen richting de oplossing}
\label{sec:Deelvragen_richting_de_oplossing}

\begin{enumerate}
  \setcounter{enumi}{3}
  \item Hoe kunnen AI-agents en workflow automatiseringsplatformen worden ingezet om informatie uit changelogs, release notes en dependency-pull-requests automatisch te interpreteren en in begrijpelijke vorm te communiceren naar het verantwoordelijke kernteam?
  \item Welke functionele en niet-functionele vereisten moet een platform voor AI-gestuurd dependencybeheer vervullen, bijvoorbeeld op het vlak van integratie met versiebeheersystemen en registries, uitbreidbaarheid, beheer van regels en policies, auditlogging en performantie?
  \item In welke mate voldoen geselecteerde AI-agent- of workflowplatformen aan deze vereisten, en hoe goed kunnen zij geïntegreerd worden in een proof-of-concept voor dependency management in de context van Springbok?
  \item In welke mate presteren de geselecteerde AI-agents effectief bij het analyseren en prioriteren van dependency-updates, gemeten in termen van nauwkeurigheid, bruikbaarheid van de adviezen en vermindering van manueel werk voor het kernteam?
\end{enumerate}

\section{\IfLanguageName{dutch}{Onderzoeksdoelstelling}{Research objective}}%
\label{sec:onderzoeksdoelstelling}

Het doel van dit onderzoek is om een proof-of-concept te ontwikkelen van een \emph{aparte webapplicatie} (dashboard) die automatisch dependency-updates detecteert, analyseert en communiceert via een geïntegreerde AI-agent, en om na te gaan in welke mate geselecteerde AI-agent- of workflowplatformen deze oplossing kunnen ondersteunen en voldoen aan de vereisten van Springbok voor geautomatiseerd dependencybeheer. De webapplicatie fungeert daarbij als losstaand platform naast GitHub en gebruikt GitHub voornamelijk als bron voor project- en dependency-informatie en als kanaal om eventuele vervolgacties (zoals pull requests of issues) te initiëren.

Concreet zal de oplossing:

\begin{itemize}
  \item per project de gebruikte dependencies bijhouden;
  \item automatisch nagaan of er nieuwe versies beschikbaar zijn via publieke registries (zoals npm);
  \item changelogs en release notes analyseren met behulp van een AI-agent die samenvat wat er gewijzigd is en welke risico's of voordelen een update inhoudt;
  \item de verantwoordelijke developers via een dashboard informeren over de aard en impact van de updates, inclusief een korte, begrijpelijke samenvatting en een indicatie of de update cruciaal, risicovol of veilig lijkt;
  \item regels ondersteunen waardoor de developers gemakkelijker kan beslissen of en wanneer een dependency geüpdatet wordt, zonder dat deze beslissingen automatisch in GitHub worden afgedwongen.
\end{itemize}

De concrete technologiekeuzes worden in de verdere uitwerking van de bachelorproef gemotiveerd op basis van de requirements en de context van Springbok. Het onderzoek resulteert in:

\begin{itemize}
  \item een prototype (demo) dat de werking van het systeem aantoont binnen een of enkele projecten van Springbok;
  \item een evaluatie van de toegevoegde waarde van AI bij dependency management voor de specifieke developer;
  \item een aanbeveling welk type platform, en eventueel welke concrete technologie, het meest geschikt is voor verdere toepassing binnen Springbok.
\end{itemize}

\section{\IfLanguageName{dutch}{Opzet van deze bachelorproef}{Structure of this bachelor thesis}}%
\label{sec:opzet-bachelorproef}

% Het is gebruikelijk aan het einde van de inleiding een overzicht te
% geven van de opbouw van de rest van de tekst. Deze sectie bevat al een aanzet
% die je kan aanvullen/aanpassen in functie van je eigen tekst.

De bachelorproef volgt een klassieke onderzoeksstructuur en is als volgt opgebouwd:

\begin{itemize}
  \item \textbf{Hoofdstuk~\ref{ch:inleiding}}
        Dit hoofdstuk beschrijft de context, de probleemstelling, de onderzoeksvragen en de doelstellingen van de bachelorproef.

  \item \textbf{Hoofdstuk~\ref{ch:stand-van-zaken}}
        In dit hoofdstuk wordt de relevante literatuur rond dependency management, bestaande tools en oplossingen, technische schuld en AI-agentframeworks en workflow-automatisering geanalyseerd.

  \item \textbf{Hoofdstuk~\ref{ch:requirementsanalyse}}
        Dit hoofdstuk bevat de requirementsanalyse voor het AI-gestuurde dependencydashboard in de context van Springbok.

  \item \textbf{Hoofdstuk~\ref{ch:methodologie}}
        Dit hoofdstuk bespreekt de onderzoeksaanpak, het ontwerp van de systeemarchitectuur, de selectie van platformen en de implementatie van de proof-of-concept.

  \item \textbf{Hoofdstuk~\ref{ch:implementatie}}
        Dit hoofdstuk licht het ontwerp en de implementatie toe van het proof-of-concept, inclusief de Mastra-agent en de n8n-workflow.

  \item \textbf{Hoofdstuk~\ref{ch:resultaten}}
        In dit hoofdstuk worden de resultaten van de implementatie en experimenten gepresenteerd, en wordt de oplossing geëvalueerd op basis van de gedefinieerde vereisten en criteria.

  \item \textbf{Hoofdstuk~\ref{ch:conclusie}}
        Dit hoofdstuk vat de belangrijkste bevindingen samen en formuleert aanbevelingen voor verdere uitwerking en mogelijke toekomstige verbeteringen.
\end{itemize}